% ====================================================================
% SECTION 05: RECOVERY OF SINGLE-MECHANISM LIMITS
% ====================================================================

\section{Recovery of Single-Mechanism Limits}

\begin{corollary}[Recovery of Single-Mechanism Limits]
Let $\eta \in [0,1]$ be the \chA{cost-mixing weight} of the unified Lagrangian
$\mathcal{L}_{\mathrm{net}}(\alpha, \eta)$. Then:
\begin{enumerate}
    \item
\textbf{Static limit} ($\eta \to 0$): The wave penalty vanishes and
$\mathcal{L}_{\mathrm{net}} \to
\mathcal{C}_{\mathrm{transport}}^{\mathrm{net}}(\alpha)$. The unique minimizer
is $\alpha^* = \alpha_t \in [\VarAlphaTLow, \VarAlphaTHigh]$, recovering the
result of Paper~I~\cite{paperI} exactly.
    \item
\textbf{Wave-dominated limit} ($\eta \to 1$): The transport penalty vanishes and
$\mathcal{L}_{\mathrm{net}} \to
\mathcal{C}_{\mathrm{wave}}^{\mathrm{net}}(\alpha)$. The unique minimizer is
$\alpha^* = \alpha_w = 2$, the impedance-matching attractor of
Paper~II~\cite{paperII}. (Note: $\alpha_w = \VarAlphaW$ is the theoretical
rigid-wall area-preserving limit for binary branching ($N=2$); the
elastic-wall-corrected value incorporating histological scaling ($h \propto
r^p$, $p = \VarP$) is $\alpha_w = \VarAlphaWFig$, as used in
Fig.~\ref{fig:phase_diagram}.)
    \item
\textbf{Minimax interior} ($\eta = \eta^*$): Neither mechanism dominates. By
Theorem~\ref{thm:arch}, the saddle point $(\alpha^*, \eta^*)$ is the unique
interior solution. The symmetric model yields $\alpha^*_{\mathrm{model}} =
\VarAlphaStarModel$, in quantitative agreement with rat pulmonary arteries
($M=\VarMassRat\,$g, $\alpha \approx \VarAlphaRatObs$, Jiang et
al.~1994~\cite{jiang1994}). The higher value observed in large mammals
($\alpha^* = \VarAlphaStar$, Kassab~1993~\cite{kassab1993}) reflects
scale-dependent morphometric heterogeneities (see \S\ref{sec:gap}).
\end{enumerate}
The unified framework therefore contains Papers~I and~II as degenerate boundary
cases of a single variational principle.
\end{corollary}

\chRB{Biologically, these two boundary limits are not abstract endpoints but
physically realized regimes---which is precisely why the single-mechanism laws
were each discovered first, in isolation. The static limit ($\eta\to0$,
$\alpha^*\to\alpha_t$) is the regime of the smallest vessels and of all
non-pulsatile transport systems: in arterioles and capillaries the Womersley
number is small, flow is effectively quasi-steady, wave reflection carries
negligible cost, and minimizing viscous dissipation alone recovers Murray's law
($\alpha\to3$)---the same exponent obeyed by the strictly non-pulsatile xylem and
airway trees. The wave-dominated limit ($\eta\to1$, $\alpha^*\to\alpha_w$) is the
regime of the largest conduit arteries, where the Womersley number is large,
pulsatile inertia dominates, and the network must instead match acoustic
impedance across each junction to avoid reflecting cardiac power back toward the
heart, recovering the area-preserving law ($\alpha\to2$). The interior minimax
$(\alpha^*,\eta^*)$ is therefore the architecture of the mid-calibre pulsatile
beds---the coronary, cerebral, and pulmonary arteries that morphometric studies
actually sample---where neither cost is negligible and the network must balance
both incommensurable penalties at once. The unified Lagrangian thus explains
\emph{why} Murray's law and the area-preserving law each hold only within their
own vascular territory, and why the universal mammalian exponent sits strictly
between them.}

\begin{remark}[Baseline vs. Elastic Wave Attractor]
The symmetric model prediction $\alpha^*_{\mathrm{model}} = \VarAlphaStarModel$
uses the baseline rigid-cylinder wave attractor $\alpha_w = 2.0$, corresponding
to acoustic impedance matching $Z \propto d^{-2}$ in the absence of wall
compliance. Paper~II~\cite{paperII} incorporates elastic wall effects via the
histological scaling $h \propto r^p$, yielding the modified wave attractor
$\alpha_w = (5-p)/2 \approx \VarAlphaWTwo$ for $p = \VarP$. The elastic
correction alone shifts the symmetric minimax to $\alpha^* \approx
\VarAlphaStarElastic$; the full coherent impedance solver of
Paper~II---integrating multiple wave reflections, bifurcation asymmetry, vessel
taper, and the Fåhræus-Lindqvist viscosity shift---recovers $\alpha^* \approx
\VarAlphaStar$ for porcine coronaries ($G = \VarG$, $N = \VarN$), in
quantitative agreement with morphometric data~\cite{kassab1993}. The present
work employs the baseline $\alpha_w = 2.0$ to isolate the fundamental
incommensurability mechanism from material-specific corrections. \chE{The}
wall thickness exponent $p = \VarP$ is an independently measured empirical
constant from Kassab's morphometric data~\cite{kassab1993}, not a free fitting
parameter used to match the predicted $\alpha^*$ to observations.
\end{remark}

Figure~\ref{fig:phase_diagram} illustrates the one-dimensional phase diagram of
the unified Lagrangian, with the two boundary attractors and the unique robust
interior saddle point.

% Sensitivity table (Theorem 2 numerical verification)

\bigskip

\noindent Theorem~\ref{thm:gauge} predicts that $\alpha^*$ is invariant to all
absolute metabolic scales and sensitive only to structural parameters.
Table~\ref{tab:sensitivity} confirms this numerically on the porcine coronary
tree ($G = \VarG$, $p = \VarP$, $\alpha_w = \VarAlphaW$, $N = \VarN$) using the
generation-by-generation network-level Lagrangian, which yields
$\alpha^*_{\mathrm{model}} = \VarAlphaStarModel$. The residual $\Delta\alpha^* =
\VarAlphaResidual$ relative to the empirically calibrated value ($\alpha^* =
\VarAlphaStar$) reflects the simplified symmetric, non-tapering architecture
used in the sensitivity model; the full morphometric model incorporating
bifurcation asymmetry and vessel taper closes the gap (see \S\ref{sec:gap}).

\begin{table}[H]
\centering
\caption{\textbf{Sensitivity analysis of $\alpha^*$ to physiological parameters.}
Log-sensitivity $S_x = \partial\alpha^*/\partial\ln x$ computed numerically on
the porcine coronary tree ($G=\VarG$, $p=\VarP$, $\alpha_w=\VarAlphaW$, $N=\VarN$)
with the network-level cost functions of Paper~II~\cite{paperII}.
\textbf{Pure metabolic parameters} (blood cost $b$, proximal flow $Q_0$, segment
length $\ell_0$) show $|S_x| < 0.01$, confirming metabolic gauge invariance
(Theorem~\ref{thm:gauge}). \chA{These are precisely the metabolic quantities
named in the abstract: the proximal flow $Q_0$ represents cardiac output, and the
blood metabolic cost $b$ encodes the ATP/oxygen stoichiometry; the table thus
directly substantiates the independence claimed there.} \textbf{Structural and fluid-mechanical parameters}
(viscosity $\mu_f$, wall metabolism $m_w$, wall exponent $p$, tree depth $G$,
wave exponent $\alpha_w$) show moderate sensitivity ($0.04 \leq |S_x| \leq 1.2$),
reflecting their role in defining the boundary conditions and topology of the
optimization problem. Despite this moderate parametric sensitivity, the predicted
$\alpha^*$ remains confined to the narrow physiological window $[\VarAlphaStarModel,
\VarAlphaStar]$
across all mammalian species.}
\label{tab:sensitivity}
\vspace{0.5em}
\begin{tabular*}{\textwidth}{@{\extracolsep{\fill}} l l l l @{}}
\toprule
Parameter ($x$) & Baseline & Perturbation & $|S_x|$ \\
\midrule
\multicolumn{4}{l}{\textit{Pure metabolic parameters (gauge-invariant: $|S_x|<0.01$)}} \\[2pt]
Blood cost ($b$)           & $\VarBblood\ \mathrm{W\,m^{-3}}$    & \VarPertB   & \VarSB   \\
Proximal flow ($Q_0$)      & $\VarQZeroML\ \mathrm{mL\,s^{-1}}$  & \VarPertQZ  & \VarSQZ  \\
Segment length ($\ell_0$)  & $\VarEllZeroMm\ \mathrm{mm}$        & \VarPertEll & \VarSEll \\
\midrule
\multicolumn{4}{l}{\textit{Physiological \& fluid-mechanical parameters (sensitive)}} \\[2pt]
Wall metabolism ($m_w$)    & $\VarMwallKW\ \mathrm{kW\,m^{-3}}$  & \VarPertMW  & \VarSMW  \\
Viscosity ($\mu_f$)        & $\VarMuFmPas\ \mathrm{mPa{\cdot}s}$ & \VarPertMuF & \VarSMuF \\
\midrule
\multicolumn{4}{l}{\textit{Structural inputs (architecture-dependent)}} \\[2pt]
Wall exponent ($p$)        & $\VarP$      & \VarPertP      & \VarSP      \\
Tree depth ($G$)           & $\VarG$      & \VarPertG      & \VarSG      \\
Wave exponent ($\alpha_w$) & $\VarAlphaW$ & \VarPertAlphaW & \VarSAlphaW \\
\bottomrule
\end{tabular*}
\end{table}

\begin{figure}[H]
\centering
\begin{tikzpicture}[>=stealth, font=\small]

  % Spine
  \draw[-] (-0.5,0) -- (11,0);

  % Axis indicators
  \draw[->] (11.3, 0.25) -- (12.3, 0.25) node[right] {$\alpha$};
  \draw[->] (12.3,-0.25) -- (11.3,-0.25) node[left]  {$\eta$};

  % Nodes
  \node[circle, fill=blue!70!black, text=white,
        minimum size=1.6cm, align=center] (W) at (1,0)
        {$\alpha_w = \VarAlphaWFig$};

  \node[circle, fill=orange!80!black, text=white,
        minimum size=2.2cm, align=center] (M) at (5.5,0)
        {$\eta^* \approx \VarEtaStar$\\[0.4ex]$\alpha^* = \VarAlphaStar$};

  \node[circle, fill=green!50!black, text=white,
        minimum size=1.6cm, align=center] (T) at (10,0)
        {$\alpha_t \approx \VarAlphaTFig$};

  % Labels below nodes
  \node[align=center, below=0.5cm of W]
        {Wave attractor\\[0.2ex]\textit{$\eta \to 1$}};
  \node[align=center, below=0.6cm of M]
        {\textbf{Minimax saddle point}};
  \node[align=center, below=0.5cm of T]
        {Transport attractor\\[0.2ex]\textit{$\eta \to 0$}};

  % Evolutionary pressure arrows
  \draw[->, dashed, bend left=30] (W.north east)
        to node[above, yshift=1mm, font=\footnotesize\itshape]
        {evolutionary pressure} (M.north west);
  \draw[->, dashed, bend right=30] (T.north west)
        to node[above, yshift=1mm, font=\footnotesize\itshape]
        {evolutionary pressure} (M.north east);

  % Bottom label
  \node[below=1.6cm of M, font=\footnotesize, text=gray]
        {unique robust interior point};

\end{tikzpicture}

\caption{\textbf{One-dimensional phase diagram of the unified network Lagrangian.}
The wave-impedance attractor ($\alpha_w = \VarAlphaWFig$,
$\eta\to 1$) and the static transport attractor ($\alpha_t \approx
\VarAlphaTFig$, $\eta\to 0$) are the two degenerate boundary cases of the
unified Lagrangian $\mathcal{L}_{\mathrm{net}}(\alpha,\eta)$. The physiological
branching exponent $\alpha^*=\VarAlphaStar$ emerges as the unique robust minimax
saddle point at \chA{marginal-balance weight} $\eta^*\approx\VarEtaStar$, stabilised by
evolutionary selection pressure from both extremes. The \chA{cost-mixing weight} $\eta$ and
the branching exponent $\alpha$ exhibit a monotonic inverse relationship at the
saddle point: increasing $\alpha$ corresponds to a shift toward metabolic
dominance (lower $\eta$).}
\label{fig:phase_diagram}
\end{figure}
